\scanpage{8}
% archive foliation: 7
% Typescript upside down in the scan; reading orientation restored here.
\begin{center}-- 37 --\end{center}
\hfill n$^{\circ}$329

assertion. On observera que $A/R_A$ s'obtient en identifiant dans $A$ les points 0 et 1. $*$

\begin{center}§ 4. Produit d'espaces topologiques.\end{center}

1. \uline{Espaces produits}

DEFINITION 1. -- \uline{Etant donnée une famille} $(X_i)_{i\in I}$ \uline{d'espaces topologiques, on appelle espace produit de cette famille l'ensemble produit} $X=\prod_{i\in I}X_i$ \uline{muni de la topologie produit des topologies des} $X_i$ (§ 2, n$^{\circ}$3, \uline{Exemple} III). \uline{On dit que les} $X_i$ $(i\in I)$ \uline{sont les espaces facteurs de} $X$.

En vertu du § 2, n$^{\circ}$ 3, prop. 4, la topologie produit sur $X$ a pour \uline{base} l'ensemble \uline{$B$} des intersections \uline{finies} d'ensembles de la forme $\operatorname{pr}_i^{-1}(U_i)$ où $U_i$ est ouvert dans $X_i$ ; ces ensembles ne sont autres que les produits $\prod_{i\in I}A_i$, où $A_i$ est \uline{ouvert} dans $X_i$ pour tout $i\in I$ et où $A_i=X_i$ \uline{sauf pour un nombre fini d'indices} ; nous donnerons à ces ensembles le nom \uline{d'ensembles élémentaires.} Si \uline{$B_i$} est une base de la topologie de $X_i$ (pour tout $i\in I$), il est clair que les ensembles élémentaires $\prod_{i\in I}A_i$ tels que $A_i\in\text{\uline{$B_i$}}$ pour tout $i$ tel que $A_i\ne X_i$ forment encore une \uline{base} de la topologie produit. Ceux de ces ensembles qui contiennent un point $x\in X$ forment par suite un \uline{système fondamental de voisinages} de $x$ (§ 1, n$^{\circ}$3, prop. 3).

Lorsque $I$ est un ensemble \uline{fini}, la construction de la topologie produit à partir des topologies des $X_i$ se simplifie : les ensembles élémentaires sont simplement les produits $\prod_{i\in I}A_i$, où $A_i$ est une partie ouverte quelconque de $X_i$ pour \uline{tout} $i\in I$ (cf. exerc.\quad).

\uline{Exemples.} -- $*$ Le produit $\unsure{\text{\uline{$R$}}|\text{\uline{$F$}}}^n$ de $n$ espaces identiques à la droite
